\chapter{Manifolds of non-negative curvature}
\label{chap:manifolds-nonneg-curvature}
\label{nonnegcurv}

\section{Busemann functions}

\begin{definition}[Minimizing geodesic ray]
\label{def:minimizing-geodesic-ray}
\lean{MorganTianLib.IsGeodesicRay}
\source{morgan-tian}
\uses{def:geodesic}
\leanok




A unit-speed geodesic ray $\lambda\colon[0,\infty)\to M$ is
\emph{minimizing} if
$$d(\lambda(s),\lambda(t))=|s-t|\qquad(s,t\ge0).$$
Equivalently, every compact subarc realizes the distance between its
endpoints. The metric identity is also used as the definition in a metric
space.
\end{definition}

\begin{definition}[Minimizing geodesics on an interval]
\label{def:minimizing-geodesic-window}
\lean{MorganTianLib.IsMinGeodesicOn}
\source{morgan-tian}
\uses{def:minimizing-geodesic-ray}
\leanok




For a metric space $(M,d)$ and $I\subset\Ar$, a map $\gamma\colon I\to M$
is a \emph{unit-speed minimizing geodesic on $I$} if
$$d(\gamma(s),\gamma(t))=|s-t|\qquad(s,t\in I).$$
The cases $I=[a,b]$, $[0,\infty)$, and $\Ar$ are called a minimizing
segment, ray, and line, respectively. Such maps are $1$-Lipschitz, and the
property is preserved by restriction and translation of the parameter.
\end{definition}

\begin{definition}[Spaces with minimizing segments]
\label{def:has-min-segments}
\lean{MorganTianLib.HasMinSegments}
\uses{def:minimizing-geodesic-window}
\leanok



A metric space $(M,d)$ \emph{has minimizing segments} if every $x,y\in M$
are joined by a unit-speed minimizing map
$\sigma\colon[0,d(x,y)]\to M$ with $\sigma(0)=x$ and
$\sigma(d(x,y))=y$. A metric space is \emph{proper} when its
closed balls are compact. Connected complete Riemannian manifolds have both
properties by Lemmas~\ref{lem:minimizing-segment-exists}
and~\ref{lem:ends-exist}.
\end{definition}

\begin{lemma}[Ultrafilter limits of minimizing geodesics]
\label{lem:ultrafilter-limit-minimizing}
\lean{MorganTianLib.exists_isMinGeodesicOn_hyperfilter_limit}
\source{morgan-tian}
\uses{def:minimizing-geodesic-window}
\leanok




Let $\cU$ be a non-principal ultrafilter on $\mathbb N$, let $M$ be a
proper metric space, and let $0\in I\subset\Ar$. Suppose $\gamma_k$ is a
unit-speed minimizing geodesic on $I_k\ni0$, every $t\in I$ belongs to
$I_k$ for cofinitely many $k$, and all $\gamma_k(0)$ lie in a compact set
$K$. Then
$$\gamma(t)=\lim_{k\to\cU}\gamma_k(t)$$
exists for every $t\in I$ and defines a unit-speed minimizing geodesic on
$I$ with $\gamma(0)\in K$.
\end{lemma}
\begin{proof}
\sketch
\leanok
\uses{def:minimizing-geodesic-window}
Choose $x_0$ and $R$ with $K\subset\overline B(x_0,R)$. For
$\cU$-almost every $k$,
$$d(\gamma_k(t),x_0)\le |t|+R,$$
so properness and compactness give the ultralimit at each $t$. Continuity
of distance yields
$$d(\gamma(s),\gamma(t))=lim_{k\to\cU}d(\gamma_k(s),\gamma_k(t))
=|s-t|.$$
Closedness of $K$ gives $\gamma(0)\in K$.
\end{proof}

\begin{lemma}[Limits of minimizing geodesics: metric version]
\label{lem:metric-limit-of-minimizing}
\lean{MorganTianLib.exists_isMinGeodesicOn_tendstoLocallyUniformlyOn}
\source{morgan-tian}
\uses{def:minimizing-geodesic-window}
\leanok




Let $M$ be proper, let $I\subset\Ar$ be a closed interval containing
$0$, and let $\gamma_k$ be unit-speed minimizing geodesics on intervals
$I_k\ni0$. Assume every point of $I$ lies in all sufficiently late
$I_k$, and $\gamma_k(0)$ lies in a fixed compact set $K$. Then a
subsequence converges uniformly on compact subsets of $I$ to a unit-speed
minimizing geodesic $\gamma\colon I\to M$ with $\gamma(0)\in K$.
\end{lemma}
\begin{proof}
\sketch
\leanok
\uses{lem:ultrafilter-limit-minimizing}
Lemma~\ref{lem:ultrafilter-limit-minimizing} gives a pointwise ultralimit
$\gamma$. On $I\cap[-n,n]$, choose a finite grid of mesh $(n+1)^{-1}$
and an increasing sequence of indices for which that interval lies in
$I_k$ and the values at every grid point differ from $\gamma$ by less
than $(n+1)^{-1}$. Since both curves are $1$-Lipschitz, interpolation
between grid points gives a uniform error at most $3/(n+1)$. This proves
local uniform convergence.
\end{proof}

\begin{lemma}[Rays from noncompact spaces: metric version]
\label{lem:ray-exists-metric}
\lean{MorganTianLib.exists_isGeodesicRay_of_noncompact}
\source{morgan-tian}
\uses{def:minimizing-geodesic-window, def:has-min-segments}
\leanok




If $M$ is a noncompact proper metric space with minimizing segments,
then every $p\in M$ is the initial point of a minimizing ray.
\end{lemma}
\begin{proof}
\sketch
\leanok
\uses{lem:ultrafilter-limit-minimizing}
Properness and noncompactness imply that $M$ is unbounded. Choose
$q_k$ with $d(p,q_k)\ge k$ and minimizing segments from $p$ to $q_k$.
After anchoring them at $p$, Lemma~\ref{lem:ultrafilter-limit-minimizing}
on $[0,\infty)$ produces the required ray.
\end{proof}

\begin{lemma}[Distinct ends differ at some compact set]
\label{lem:ends-differ-at-compact}
\lean{MorganTianLib.SpaceOfEnds.exists_apply_ne}
\uses{def:space-of-ends}
\leanok



If $\cE_1\ne\cE_2$ are ends of a topological space $X$, then some
compact $K\subset X$ determines different components of $X\setminus K$
for $\cE_1$ and $\cE_2$.
\end{lemma}
\begin{proof}
\sketch
\leanok
Ends are compatible families in the inverse limit
$\varprojlim_K\pi_0(X\setminus K)$; unequal families differ at some
coordinate.
\end{proof}

\begin{lemma}[Ends reach infinity]
\label{lem:end-component-unbounded}
\lean{MorganTianLib.SpaceOfEnds.exists_mem_connectedComponentIn_forall_le_dist}
\source{morgan-tian}
\uses{def:space-of-ends, def:ends-transition}
\leanok




Let $M$ be proper, let $\cE$ be an end, and let $U$ be the component of
$M\setminus K$ selected by $\cE$ for a compact $K$. For every finite
$R$ there is $z\in U$ such that $d(z,w)\ge R$ for every $w\in K$.
\end{lemma}
\begin{proof}
\sketch
\leanok
\uses{def:ends-transition}
Otherwise $U$ is bounded, hence $\overline U$ is compact. For
$K'=K\cup\overline U$, compatibility of the end under the transition
map forces the selected nonempty component of $M\setminus K'$ to lie in
$U\subset K'$, a contradiction.
\end{proof}

\begin{lemma}[Lines from two ends: metric version]
\label{lem:line-exists-metric}
\lean{MorganTianLib.exists_isMinGeodesicOn_univ_of_ends_ne}
\source{morgan-tian}
\uses{def:minimizing-geodesic-window, def:has-min-segments, def:space-of-ends}
\leanok




A proper metric space with minimizing segments and at least two ends
contains a minimizing line.
\end{lemma}
\begin{proof}
\sketch
\leanok
\uses{lem:ultrafilter-limit-minimizing,lem:ends-differ-at-compact,lem:end-component-unbounded}
Choose a compact $K$ separating components $U_1,U_2$ selected by two
ends. Pick $x_k\in U_1$ and $y_k\in U_2$ at distance at least $k$ from
$K$. Every minimizing segment from $x_k$ to $y_k$ meets $K$; recenter at
such an intersection. Both parameter endpoints then tend to infinity,
while the new origins remain in $K$. The ultrafilter limit lemma with
$I=\Ar$ gives a minimizing line.
\end{proof}

\begin{lemma}[Geodesics depend continuously on their initial conditions]
\label{lem:geodesic-continuous-dependence}
\lean{MorganTianLib.exists_globalGeodesic_initial, MorganTianLib.eqOn_of_initial_data, MorganTianLib.convAt_of_convAt_zero, MorganTianLib.tendsto_apply_of_convAt_zero, MorganTianLib.tendstoUniformlyOn_of_convAt_zero, MorganTianLib.convAt_zero_of_tendsto_tangentBundle, MorganTianLib.tendsto_tangentBundle_velocity_of_convAt, MorganTianLib.globalGeodesic_initial_state, MorganTianLib.globalGeodesic_tangentBundle_continuousDependence}
\source{morgan-tian}
\uses{def:geodesic, thm:hopf-rinow}

\leanok




Let $(M,g)$ be complete.
\begin{enumerate}
\item Every $(q,v)\in TM$ determines a unique global geodesic
$\gamma_{q,v}$ with initial data $(q,v)$, and every geodesic on an
interval containing $0$ with those data is its restriction.
\item If $(q_k,v_k)\to(q,v)$ in $TM$, then for each $T<\infty$,
$$\sup_{|t|\le T}d(\gamma_{q_k,v_k}(t),\gamma_{q,v}(t))\to0,$$
and the position-velocity pairs converge in $TM$ at each $t\in\Ar$.
\end{enumerate}
\end{lemma}
\begin{proof}
\sketch
\uses{def:geodesic,thm:hopf-rinow}
In a chart the geodesic state $Z=(x,\dot x)$ solves
$$\dot Z=F(Z),\qquad F(x,y)=(y,-\Gamma^k_{ij}(x)y^iy^j).$$
On a compact coordinate ball, $F$ is bounded and Lipschitz. A no-escape
estimate keeps nearby solutions in the ball for a uniform time, and
Gr\"onwall gives
$$|Z_1(t)-Z_2(t)|\le e^{L|t-t_0|}|Z_1(t_0)-Z_2(t_0)|.$$
Local existence and Hopf--Rinow give global existence; the estimate and
a connectedness argument give uniqueness. Applying the same estimate on
successive overlapping coordinate intervals along $\gamma_{q,v}$ proves
state convergence and uniform convergence of positions on every compact
time interval.
\end{proof}

\begin{lemma}[Limits of minimizing geodesics]
\label{lem:limit-of-minimizing-geodesics}
\lean{MorganTianLib.exists_isMinGeodesicOn_convAt_of_tendsto}
\source{morgan-tian}
\uses{def:geodesic}




Let $M$ be complete, let $I\subset\Ar$ be a nondegenerate interval
containing $0$, and let closed intervals $I_k\ni0$ eventually contain
every compact subset of $I$. If unit-speed minimizing geodesics
$\gamma_k\colon I_k\to M$ satisfy $\gamma_k(0)\to p$, then a subsequence
converges uniformly on compact subsets of $I$ to a unit-speed minimizing
geodesic $\gamma\colon I\to M$ with $\gamma(0)=p$, and its initial
position-velocity pairs converge in $TM$.
\end{lemma}
\begin{proof}
\sketch
\uses{lem:geodesic-continuous-dependence,thm:hopf-rinow,thm:levi-civita-connection,def:geodesic}
Extend each $\gamma_k$ globally. In a fixed chart near $p$, unit speed
uniformly bounds the coordinate components of $\gamma_k'(0)$; extract a
subsequence converging to a unit $v\in T_pM$. Continuous dependence gives
local uniform convergence to $\gamma_{p,v}$ on $I$. Passing to the limit
in
$d(\gamma_k(s),\gamma_k(t))=|s-t|$ proves that the limit is minimizing.
\end{proof}

\begin{lemma}[Minimizing segments exist in complete manifolds]
\label{lem:minimizing-segment-exists}
\lean{MorganTianLib.hasMinSegments_of_complete}
\uses{def:geodesic}
\leanok



If $M$ is connected and complete, any distinct $x,q\in M$ are joined by
a unit-speed minimizing geodesic on $[0,d(x,q)]$.
\end{lemma}
\begin{proof}
\sketch
\leanok
\uses{def:cut-locus,lem:cut-time-star-shaped,def:exponential-map,thm:hopf-rinow,lem:limit-of-minimizing-geodesics}
Outside the cut locus the assertion follows from the star-shaped
minimizing domain of $\exp_x$. Approximate a point in the cut locus by
points $q_k$ outside it and take the limit of their minimizing segments
using Lemma~\ref{lem:limit-of-minimizing-geodesics}. Completeness extends
the limiting geodesic to $d(x,q)$, and the endpoint estimate
$d(\sigma(s),q)\le d(x,q)-s$ forces $\sigma(d(x,q))=q$.
\end{proof}

\begin{lemma}[Non-compact manifolds have ends]
\label{lem:ends-exist}
\lean{MorganTianLib.exists_isGeodesicRay_of_complete_noncompact', MorganTianLib.exists_isMinGeodesicOn_univ_of_complete_ends_ne', MorganTianLib.properSpace_of_complete', MorganTianLib.geodesicEndpointContinuity, MorganTianLib.convStepProperty}
\source{morgan-tian}
\uses{def:geodesic, def:space-of-ends}
\leanok
\label{ends}




If $M$ is complete, connected, and noncompact, then every $p\in M$ is
the initial point of a minimizing ray. If $M$ has more than one end, it
also contains a minimizing line.
\end{lemma}
\begin{proof}
\sketch
\leanok
\uses{def:geodesic,thm:hopf-rinow,def:exponential-map,lem:minimizing-segment-exists,def:has-min-segments,lem:ray-exists-metric,lem:line-exists-metric,def:space-of-ends}
Lemma~\ref{lem:minimizing-segment-exists} gives minimizing segments.
Hopf--Rinow and the exponential map identify each closed ball with the
image of the compact tangent ball of the same radius, so $M$ is proper.
Apply Lemma~\ref{lem:ray-exists-metric}, and, when there are two ends,
Lemma~\ref{lem:line-exists-metric}.
\end{proof}

\begin{definition}[Busemann approximants]
\label{def:busemann-approx}
\lean{MorganTianLib.busemannAux}
\source{morgan-tian}
\uses{def:minimizing-geodesic-ray}
\leanok




For a minimizing ray $\lambda$ and $t\ge0$, define
$$B_{\lambda,t}(x)=d(\lambda(t),x)-t.$$
\end{definition}

\begin{lemma}[Busemann approximants are $1$-Lipschitz]
\label{lem:busemann-approx-lipschitz}
\lean{MorganTianLib.lipschitzWith_busemannAux}
\uses{def:busemann-approx}
\leanok



For all $x,y\in M$ and $t\ge0$,
$|B_{\lambda,t}(x)-B_{\lambda,t}(y)|\le d(x,y)$.
\label{eqnBlip}
\end{lemma}
\begin{proof}
\sketch
\leanok
This is the reverse triangle inequality for distance from $\lambda(t)$.
\end{proof}

\begin{lemma}[Busemann approximants are non-increasing in $t$]
\label{lem:busemann-approx-antitone}
\lean{MorganTianLib.busemannAux_antitoneOn}
\uses{def:busemann-approx}
\leanok



For each $x$, the function $t\mapsto B_{\lambda,t}(x)$ is non-increasing
on $[0,\infty)$.
\label{eqnBmono}
\end{lemma}
\begin{proof}
\sketch
\leanok
For $s\le t$, the triangle inequality and minimality give
$d(\lambda(t),x)\le(t-s)+d(\lambda(s),x)$.
\end{proof}

\begin{lemma}[Busemann approximants are bounded below]
\label{lem:busemann-approx-lower-bound}
\lean{MorganTianLib.neg_dist_le_busemannAux}
\uses{def:busemann-approx}



For $x\in M$ and $t\ge0$,
$$B_{\lambda,t}(x)\ge-d(x,\lambda(0));$$
in particular $B_{\lambda,t}(\lambda(0))=0$.
\end{lemma}
\begin{proof}
\sketch
\leanok
Use $t=d(\lambda(0),\lambda(t))\le d(\lambda(0),x)+d(x,\lambda(t))$.
\end{proof}

\begin{definition}[Busemann function]
\label{def:busemann-function}
\lean{MorganTianLib.busemann, MorganTianLib.busemann_spec_of_ray}
\source{morgan-tian}
\uses{def:busemann-approx, lem:busemann-approx-antitone, lem:busemann-approx-lower-bound, lem:ends-exist}
\leanok




The \emph{Busemann function} of a minimizing ray $\lambda$ is
$$B_\lambda(x)=\lim_{t\to\infty}B_{\lambda,t}(x)
=\inf_{t\ge0}B_{\lambda,t}(x),$$
whose existence follows from monotonicity and the preceding lower bound.
\end{definition}

\begin{lemma}[Busemann approximants converge to the Busemann function]
\label{lem:busemann-limit}
\lean{MorganTianLib.tendsto_busemannAux_atTop}
\uses{def:busemann-function}
\leanok



For each $x\in M$, $B_{\lambda,t}(x)\to B_\lambda(x)$ as
$t\to\infty$.
\end{lemma}
\begin{proof}
\sketch
\leanok
\uses{lem:busemann-approx-antitone,lem:busemann-approx-lower-bound}
A decreasing real function bounded below converges to its infimum.
\end{proof}

\begin{lemma}[The Busemann function bounds its approximants from below]
\label{lem:busemann-le-approx}
\lean{MorganTianLib.busemann_le_busemannAux}
\uses{def:busemann-function}
\leanok



$B_\lambda(y)\le B_{\lambda,t}(y)$ for every $y\in M$ and $t\ge0$.
\end{lemma}
\begin{proof}
\sketch
\leanok
The left side is the infimum over all $t\ge0$.
\end{proof}

\begin{lemma}[The Busemann function is $1$-Lipschitz]
\label{lem:busemann-lipschitz}
\lean{MorganTianLib.lipschitzWith_busemann}
\uses{def:busemann-function}
\leanok



For all $x,y\in M$,
$$|B_\lambda(x)-B_\lambda(y)|\le d(x,y).$$
Thus $B_\lambda$ is continuous.
\end{lemma}
\begin{proof}
\sketch
\leanok
\uses{lem:busemann-approx-lipschitz}
Pass the $1$-Lipschitz inequalities for $B_{\lambda,t}$ to their
pointwise limit.
\end{proof}

\begin{lemma}[Lower bound for the Busemann function]
\label{lem:busemann-lower-bound}
\lean{MorganTianLib.neg_dist_le_busemann}
\uses{def:busemann-function}
\leanok



$B_\lambda(x)\ge-d(x,\lambda(0))$ for all $x\in M$.
\end{lemma}
\begin{proof}
\sketch
\leanok
\uses{lem:busemann-approx-lower-bound}
Take the infimum in Lemma~\ref{lem:busemann-approx-lower-bound}.
\end{proof}

\begin{lemma}[The Busemann function along its ray]
\label{lem:busemann-along-ray}
\lean{MorganTianLib.busemann_apply_ray}
\uses{def:busemann-function}
\leanok



$B_\lambda(\lambda(s))=-s$ for all $s\ge0$.
\end{lemma}
\begin{proof}
\sketch
\leanok
$B_{\lambda,t}(\lambda(s))=|s-t|-t$, whose infimum over $t\ge0$ is
$-s$.
\end{proof}

\begin{definition}[Asymptotic rays]
\label{def:asymptotic-ray}
\lean{MorganTianLib.IsAsymptoticRay}
\source{morgan-tian}
\uses{def:minimizing-geodesic-ray, def:minimizing-geodesic-window}
\leanok




Let $\lambda$ and $\mu$ be minimizing rays in a metric space. The ray
$\mu$ is \emph{asymptotic to $\lambda$} if there are $t_n\to\infty$
and minimizing segments from $\mu(0)$ to $\lambda(t_n)$ which converge
to $\mu$ uniformly on compact subsets of $[0,\infty)$. Their domains
exhaust $[0,\infty)$ because
$d(\mu(0),\lambda(t_n))\ge t_n-d(\mu(0),\lambda(0))$.
\end{definition}

\begin{lemma}[Existence of asymptotic rays]
\label{lem:asymptotic-ray-exists}
\lean{MorganTianLib.exists_isAsymptoticRay}
\source{morgan-tian}
\uses{def:asymptotic-ray, def:has-min-segments}
\leanok




In a proper metric space with minimizing segments, every $x$ is the
initial point of a ray asymptotic to any prescribed minimizing ray
$\lambda$.
\end{lemma}
\begin{proof}
\sketch
\leanok
\uses{lem:metric-limit-of-minimizing}
Take minimizing segments from $x$ to $\lambda(n)$. Their lengths tend
to infinity, and Lemma~\ref{lem:metric-limit-of-minimizing} supplies a
locally uniform subsequential limit.
\end{proof}

\begin{lemma}[Busemann values along asymptotic rays]
\label{lem:asymptotic-ray}
\lean{MorganTianLib.IsAsymptoticRay.busemann_apply}
\uses{def:busemann-function, def:asymptotic-ray}
\leanok



If $\mu$ is asymptotic to $\lambda$ and $x=\mu(0)$, then
$$B_\lambda(\mu(s))=B_\lambda(x)-s\qquad(s\ge0).$$
\end{lemma}
\begin{proof}
\sketch
\leanok
\uses{lem:busemann-approx-antitone,lem:busemann-limit,lem:busemann-lipschitz,lem:busemann-le-approx,lem:busemann-along-ray}
For witnessing segments $\sigma_n$ from $x$ to $\lambda(t_n)$,
minimality gives
$$B_{\lambda,t_n}(\sigma_n(s))=B_{\lambda,t_n}(x)-s.$$
Use local uniform convergence, monotonicity of the approximants, and
their convergence at $x$ to obtain
$B_\lambda(\mu(s))\le B_\lambda(x)-s$. The reverse inequality follows
from the $1$-Lipschitz property and $d(x,\mu(s))=s$.
\end{proof}

\begin{corollary}[Busemann upper bound from an asymptotic ray]
\label{cor:asymptotic-ray-busemann-bound}
\lean{MorganTianLib.IsAsymptoticRay.busemann_le}
\uses{def:busemann-function, def:asymptotic-ray, lem:asymptotic-ray}
\leanok



If $x=\mu(0)$ and $\mu$ is asymptotic to $\lambda$, then
$$B_\lambda(y)\le B_\lambda(x)-t+d(y,\mu(t))$$
for every $y\in M$ and $t\ge0$.
\end{corollary}
\begin{proof}
\sketch
\leanok
\uses{lem:busemann-lipschitz,lem:asymptotic-ray}
Apply the Lipschitz bound between $y$ and $\mu(t)$ and then use
Lemma~\ref{lem:asymptotic-ray}.
\end{proof}

\begin{lemma}[Busemann gradient along an asymptotic ray]
\label{lem:busemann-gradient-norm-one}
\uses{def:busemann-function, def:asymptotic-ray, lem:asymptotic-ray-exists, lem:asymptotic-ray}

Let $M$ be connected and complete, let $\mu$ be asymptotic to $\lambda$
with $\mu(0)=x$, and suppose $B_\lambda$ is differentiable at $x$. Then
$$\nabla B_\lambda(x)=-\mu'(0),\qquad |\nabla B_\lambda(x)|=1.$$
\label{eqnBlipgrad}
\end{lemma}
\begin{proof}
\sketch
\uses{def:geodesic,thm:levi-civita-connection}
Differentiating $B_\lambda(\mu(s))=B_\lambda(x)-s$ at $s=0$ gives
$\langle\nabla B_\lambda(x),\mu'(0)\rangle=-1$. Differentiability and
the $1$-Lipschitz bound imply
$|\langle\nabla B_\lambda(x),v\rangle|\le1$ for every unit $v$, by
testing along unit-speed geodesics. Equality in Cauchy--Schwarz then
forces $\nabla B_\lambda(x)=-\mu'(0)$.
\end{proof}

\begin{proposition}[Busemann function has non-positive Laplacian]
\label{prop:busemann-laplacian}
\source{morgan-tian}
\uses{def:busemann-function, def:ricci-curvature}
\label{Blambda}


If $M$ is complete with non-negative Ricci curvature, then every
Busemann function satisfies $\Delta B_\lambda\le0$ in the weak sense.
\end{proposition}
\begin{proof}
\sketch
\uses{def:busemann-function,ex:distance-function-laplacian,rem:laplacian-weak-sense,thm:hopf-rinow,def:geodesic,lem:laplacian-local-formula}
This is the Schoen--Yau argument \cite{SchoenYau}.
For $n=1$, completeness and the existence of the ray identify $M$
isometrically with $\Ar$; in the induced coordinate
$B_\lambda(s)=-s$, so its weak Laplacian vanishes. Assume $n\ge2$.
For a non-negative test function $\varphi$ supported in $C$ and
$f_t(x)=d(\lambda(t),x)$, Laplacian comparison gives
$$\int_M B_{\lambda,t}\Delta\varphi\,d\vol
\le\int_M\frac{n-1}{f_t}\varphi\,d\vol.$$
Writing $p=\lambda(0)$ and
$D=\max_{x\in C}d(p,x)$, one has $f_t\ge t-D$ on $C$, hence the
right side is at most
$(n-1)(t-D)^{-1}\int_M\varphi\,d\vol$ for $t>D$. The approximants
converge pointwise to $B_\lambda$ and are uniformly bounded on $C$, so
dominated convergence passes to the limit on the left. Letting
$t\to\infty$ yields
$\int_M B_\lambda\Delta\varphi\,d\vol\le0$, which is the stated weak
inequality.
\end{proof}

\section{Comparison results in non-negative curvature}
\label{eqngradvanish}
\label{eqnradial}
\label{locallabel}
\label{eq:tangential-acceleration}
\label{eq:first-variation-energy}
\label{eq:laplacian-frame-sum}
\label{eq:dual-metric-compat}

These are the Toponogov comparisons \cite{Toponogov}; proofs may be
found in \cite{Sakai} and \cite{CheegerEbin}; see also \cite{Petersen}.

For a triangle with minimizing sides $s_{ab},s_{ac},s_{bc}$, write
$|s_{xy}|=d(x,y)$ and let $\angle_a$ be the angle between $s_{ab}$ and
$s_{ac}$ at $a$.

\begin{theorem}[Length comparison theorem]
\label{thm:length-comparison}
\uses{def:sectional-curvature}
\label{lengthcompar}

Let $M$ have non-negative sectional curvature, let
$\triangle(a,b,c)$ be a geodesic triangle in $M$, and let
$\triangle(a',b',c')$ be Euclidean.
\begin{enumerate}
\item If corresponding side lengths agree, each Euclidean angle is at
most its corresponding angle in $M$. Moreover, if $x,y$ lie at distances
$\alpha,\beta$ from $a$ along $s_{ab},s_{ac}$ and $x',y'$ are the
corresponding Euclidean points, then $d(x,y)\ge d(x',y')$.
\item If the two sides adjacent to $a,a'$ and the included angles agree,
then $|s_{b'c'}|\ge|s_{bc}|$.
\end{enumerate}
\end{theorem}

For $u\in[0,|s_{ab}|]$ and $v\in[0,|s_{ac}|]$, let $EA(u,v)$ be the
angle opposite $d(x(u),y(v))$ in the Euclidean triangle with adjacent
sides $u,v$, where $x(u)\in s_{ab}$ and $y(v)\in s_{ac}$ are at
distances $u,v$ from $a$.

\begin{corollary}[Angle monotonicity in non-negative curvature]
\label{cor:angle-monotonicity}
\uses{thm:length-comparison}
\label{anglemono}

The function $EA(u,v)$ is non-increasing in either variable when the
other is fixed.
\end{corollary}

\begin{corollary}[Angle comparison in non-negative curvature]
\label{cor:angle-comparison}
\uses{thm:length-comparison}
\label{anglecompar}

Let $M$ be complete with non-negative curvature. For minimizing arcs
from $p$ to $a,b,c$, complete the three geodesic triangles
$T(a,p,b)$, $T(b,p,c)$, and $T(c,p,a)$, and let the corresponding
Euclidean comparison triangles have vertices $a',p',b',c'$. Then
$$\angle_{p'}T(a',p',b')+\angle_{p'}T(b',p',c')
+\angle_{p'}T(c',p',a')\le2\pi.$$
\end{corollary}
\begin{proof}
\sketch
\uses{thm:length-comparison}
Shorten the three radial sides while fixing their directions. Length
comparison makes the comparison angles non-increasing; their limiting
sum at $p$ is the sum of three Euclidean angles between tangent vectors,
which is at most $2\pi$.
\end{proof}

\section{The soul theorem}

A subset is \emph{totally convex} if every minimizing geodesic segment
with endpoints in it remains in it.

\begin{theorem}[Soul Theorem (Cheeger-Gromoll)]
\label{thm:soul-theorem}
\uses{def:sectional-curvature}
\label{soul}

Let $(M,g)$ be connected, complete, noncompact, and of non-negative
sectional curvature. Then $M$ contains a compact, totally geodesic,
totally convex submanifold $S$ of positive codimension and is
diffeomorphic to the total space of its normal bundle. If the curvature
is positive, every such soul is a point and $M\cong\Ar^n$.
This is the Cheeger--Gromoll soul theorem
\cite{CheegerGromollbul,CheegerGromoll}.
\end{theorem}

\begin{remark}[Soul theorem for positive curvature]
\label{rem:soul-theorem-positive-curvature}
\uses{thm:soul-theorem}

The positive-curvature case used below has a point soul; this case first
appeared in \cite{GromollMeyer}.
\end{remark}

The remaining soul-theorem consequences follow the comparison-geometric
construction in \cite{Petersen}.

\begin{lemma}[Angles of geodesics converge at infinity]
\label{lem:angle-convergence-infinity}
\uses{def:sectional-curvature}

Let $M$ be complete, noncompact, and of non-negative sectional
curvature. For $p\in M$ and $\epsilon>0$, some compact $K$ has the
property that any two minimizing geodesics from $p$ to a common
$q\notin K$ meet at $q$ at angle less than $\epsilon$.
\end{lemma}
\begin{proof}
\sketch
\uses{thm:length-comparison}
Otherwise choose $q_n\to\infty$ and pairs $\gamma_n,\mu_n$ whose angle
at $q_n$ is at least $\epsilon$. Pass to a subsequence so the initial
directions within each family differ by at most $O(\epsilon)$ and
$d(p,q_{n+1})\ge100\epsilon^{-2}d(p,q_n)$. Applying length comparison
to triangles $(p,q_n,q_{n+1})$ shows that both terminal directions at
$q_n$ make angle less than $\epsilon/2$ with the reverse direction of
a minimizing segment to $q_{n+1}$. Their mutual angle is then less than
$\epsilon$, a contradiction.
\end{proof}

\begin{corollary}[Distance level sets are homeomorphic]
\label{cor:distance-level-sets-homeomorphic}
\uses{lem:angle-convergence-infinity}
\label{homeo}

Let $M$ be complete, noncompact, and of non-negative sectional
curvature, fix $p\in M$, and put $f=d(p,\cdot)$. For some $R<\infty$
and all $R\le s<s'$,
\begin{enumerate}
\item $f^{-1}([s,s'])\cong f^{-1}(s)\times[s,s']$;
\item $f^{-1}([s,\infty))\cong f^{-1}(s)\times[s,\infty)$.
\end{enumerate}
\end{corollary}
\begin{proof}
\sketch
The outward field construction is given in \cite{Petersen}.
Choose $R$ so that minimizing directions from $p$ to points outside
$B(p,R/2)$ lie in cones of angle at most $\pi/6$. A smooth unit vector
field in these cones satisfies $Xf\ge\cos(\pi/3)>0$. Each flow line
therefore crosses every level $f^{-1}(s)$, $s\ge R$, exactly once, and
the flow gives both product homeomorphisms.
\end{proof}

In positive curvature the soul is a point. The outward vector-field
construction may be chosen on all of $M$, stationary only at the soul;
it identifies every positive distance level with $S^{n-1}$ and every
annulus between such levels with a product \cite{Petersen}.

\begin{corollary}[Distance annuli are homotopy spheres]
\label{cor:distance-annuli-homotopy-sphere}
\uses{cor:distance-level-sets-homeomorphic, thm:soul-theorem}
\label{htpytype}

If $M$ is a complete noncompact positively curved $n$-manifold and
$f=d(p,\cdot)$, then for some $R=R(p)$ and every $R\le s<s'$, both
$f^{-1}((s,s'))$ and $f^{-1}((s,\infty))$ are homotopy equivalent to
$S^{n-1}$.
\end{corollary}
\begin{proof}
\sketch
\uses{cor:distance-level-sets-homeomorphic}
The soul theorem gives $M\cong\Ar^n$. Far enough out, the product
description of Corollary~\ref{cor:distance-level-sets-homeomorphic}
makes all tails homotopy equivalent. Exhaustion by balls shows that a
tail is simply connected and has trivial $H_i$ for $i<n-1$; noncompact
Poincar\'e duality gives $H_{n-1}\cong\Zee$. Hurewicz identifies its
homotopy type with $S^{n-1}$, and the annulus has the same type by the
product description.
\end{proof}

\section{Ends of a manifold}
\label{ENDS}

\begin{definition}[Transition maps between complementary components]
\label{def:ends-transition}
\lean{MorganTianLib.endsTransition}
\source{morgan-tian}
\leanok



For compact $K\subset K'\subset X$, inclusion of complements induces
$$\pi_0(X\setminus K')\longrightarrow\pi_0(X\setminus K),$$
sending each component to the unique component containing it. Each
$\pi_0$ carries the quotient topology.
\end{definition}

\begin{lemma}[The transition maps are continuous]
\label{lem:ends-transition-continuous}
\lean{MorganTianLib.continuous_endsTransition}
\uses{def:ends-transition}
\leanok



Every transition map of Definition~\ref{def:ends-transition} is
continuous.
\end{lemma}
\begin{proof}
\sketch
\leanok
The continuous map $X\setminus K'\to\pi_0(X\setminus K)$ is constant
on components and therefore descends through the quotient.
\end{proof}

\begin{lemma}[Transition maps: identity]
\label{lem:ends-transition-id}
\lean{MorganTianLib.endsTransition_refl}
\uses{def:ends-transition}
\leanok



The transition map associated with $K\subset K$ is the identity.
\end{lemma}
\begin{proof}
\sketch
\leanok
Each component is the unique component containing itself.
\end{proof}

\begin{lemma}[Transition maps: composition]
\label{lem:ends-transition-comp}
\lean{MorganTianLib.endsTransition_trans}
\uses{def:ends-transition}
\leanok



For compact $K\subset K'\subset K''$, the transition from $K''$ to
$K$ is the composite through $K'$. Thus $K\mapsto\pi_0(X\setminus K)$
is an inverse system.
\end{lemma}
\begin{proof}
\sketch
\leanok
This follows from transitivity and uniqueness of the containing
component.
\end{proof}

\begin{definition}[Space of ends]
\label{def:space-of-ends}
\lean{MorganTianLib.SpaceOfEnds}
\source{morgan-tian}
\uses{def:ends-transition, lem:ends-transition-comp, lem:compact-set-in-codim-zero-submanifold}





For a connected manifold $M$, its \emph{space of ends} is
$$\varprojlim_{K\subset M\text{ compact}}\pi_0(M\setminus K),$$
with the subspace topology from the product. An end is therefore a
compatible choice of complementary component for every compact $K$.
Compact codimension-zero submanifolds are cofinal among compact sets by
Lemma~\ref{lem:compact-set-in-codim-zero-submanifold}, yielding the
classical equivalent indexing; their complementary component sets are
finite. Since complements are locally connected, their components are
open and their quotient component spaces are discrete. The selected
components are neighborhoods of the end, and a sequence converges to an
end when it is eventually in each of them.
\end{definition}

\begin{lemma}[Compact sets lie in compact codimension-zero submanifolds]
\label{lem:compact-set-in-codim-zero-submanifold}
Every compact subset $C$ of a smooth $n$-manifold lies in a compact
codimension-zero smooth submanifold with boundary.
\end{lemma}
\begin{proof}
\sketch
Choose finitely many compactly supported coordinate bump functions whose
positive sets cover $C$, and let $\psi$ be their sum. For
$0<c<\min_C\psi$ chosen as a regular value, the superlevel set
$K=\psi^{-1}([c,\infty))$ is compact, contains $C$, and has smooth
boundary $\psi^{-1}(c)$ by the regular-value theorem.
\end{proof}

A proper path $\gamma\colon[a,b)\to M$ tends to an end $\cE$ if it is
eventually in every component selected by $\cE$. An end of a Riemannian
manifold is at \emph{finite distance} when such a path can have finite
length; the distance from a point to that end is the infimum of these
lengths, and is positive. A Riemannian manifold is complete exactly when it has no
finite-distance end.


\section{The splitting theorem}

The splitting theorem is due to Cheeger--Gromoll
\cite{CheegerGromoll2}; Toponogov proved the sectional-curvature version
\cite{Toponogov}.

\begin{lemma}[Green identity for compactly supported test functions]
\label{lem:green-identity-compact-support}
\lean{MorganTianLib.green_identity_compact_support}
\uses{lem:laplacian-local-formula}
\leanok

If $u$ is smooth on an open $\Omega\subset M$ and
$\psi\in C_c^\infty(\Omega)$, then
$$\int_\Omega(u\Delta\psi-\psi\Delta u)\,d\vol=0.$$
\end{lemma}
\begin{proof}
\sketch
\uses{lem:laplacian-local-formula}
In one coordinate chart the integrand times the volume density is
$$\partial_i\!\left(g^{ij}\sqrt{\det g},
(u\partial_j\psi-\psi\partial_ju)\right),$$
whose integral vanishes by compact support. A partition of unity reduces
the general case to this calculation.
\end{proof}

\begin{claim}[Dirichlet problem on small geodesic balls]
\label{claim:dirichlet-problem-small-balls}
\uses{def:exponential-map, cor:exponential-local-diffeomorphism, lem:laplacian-local-formula}

For every $p$ in an $n$-dimensional Riemannian manifold there is
$r_D(p)>0$ such that, for $0<s\le r_D(p)$:
\begin{enumerate}
\item $\exp_p$ identifies $B(0,s)\subset T_pM$ diffeomorphically with
$B(p,s)$, identifies their closures, and identifies the boundary with
the smooth metric sphere $S(p,s)\cong S^{n-1}$;
\item every continuous $\varphi\colon S(p,s)\to\Ar$ has a continuous
extension $h$ to $\overline B(p,s)$ which is smooth and harmonic in
$B(p,s)$.
\end{enumerate}
The normal-ball assertion follows from the exponential-map theory in
\cite{doCarmo}; the Dirichlet assertion is the standard local elliptic
result used in \cite{Petersen}.
\end{claim}

\begin{claim}[Comparison principle for continuous weak subsolutions]
\label{claim:weak-subsolution-comparison}
\uses{rem:laplacian-weak-sense, lem:laplacian-local-formula, lem:green-identity-compact-support}

Let $\Omega$ be a nonempty connected open subset of a Riemannian
manifold, with compact closure and nonempty boundary. If
$w\in C(\overline\Omega)$ satisfies
$$\int_\Omega w\Delta\psi\,d\vol\ge0$$
for every non-negative $\psi\in C_c^\infty(\Omega)$, then
$$\max_{\overline\Omega}w=\max_{\partial\Omega}w.$$
This is the distributional comparison principle; it is distinct from
the barrier formulation used in \cite{Petersen}.
\end{claim}

\begin{lemma}[Second-derivative test, minimum version]
\label{lem:second-derivative-test-min}
\mathlibok
Let $\phi$ be continuous at $t_0$, differentiable near $t_0$, with
$\phi'(t_0)=0$ and with $\phi'$ differentiable at $t_0$. If
$\phi''(t_0)>0$, then $t_0$ is a local minimum.
\end{lemma}

\begin{lemma}[Second-derivative test at a local maximum]
\label{lem:second-derivative-test-max}
\lean{MorganTianLib.deriv_deriv_nonpos_of_isLocalMax}
\leanok


If $\phi$ is continuous and locally maximal at $t_0$, differentiable
near $t_0$, and $\phi'$ is differentiable at $t_0$, then
$\phi''(t_0)\le0$.
\end{lemma}
\begin{proof}
\sketch
\leanok
\uses{lem:second-derivative-test-min}
Fermat's lemma gives $\phi'(t_0)=0$. A positive second derivative would
also make $t_0$ a local minimum; local constancy would then force
$\phi''(t_0)=0$.
\end{proof}

\begin{lemma}[Existence of integral curves]
\label{lem:integral-curve-exists}
\mathlibok
For a smooth vector field $X$ on a boundaryless smooth manifold and
$q\in M$, there is a curve $\gamma\colon\Ar\to M$ with $\gamma(0)=q$
which satisfies $\gamma'=X\circ\gamma$ near $0$.
\end{lemma}

\begin{lemma}[Vanishing differential at a local maximum]
\label{lem:differential-vanishes-at-max}
\lean{MorganTianLib.mfderiv_eq_zero_of_isLocalMax}


If a smooth $f\colon M\to\Ar$ has a local maximum at $q$, then
$df_q=0$.
\end{lemma}
\begin{proof}
\sketch
\leanok
\uses{lem:integral-curve-exists}
Extend any $v\in T_qM$ to a smooth field and evaluate $f$ along its
integral curve through $q$. Fermat's lemma gives $df_q(v)=0$.
\end{proof}

\begin{lemma}[Hessian is negative semi-definite at a local maximum]
\label{lem:hessian-nonpositive-at-max}
\lean{MorganTianLib.hessianAt_nonpos_of_isLocalMax}
\uses{lem:hessian-symmetric}



For a smooth function with a local maximum at $q$ and any affine
connection,
$$\Hess(f)_q(v,v)\le0\qquad(v\in T_qM).$$
\end{lemma}
\begin{proof}
\sketch
\leanok
\uses{lem:differential-vanishes-at-max,lem:integral-curve-exists,lem:second-derivative-test-max,lem:hessian-symmetric}
Extend $v$ to a field $X$ and let $\gamma$ be its integral curve.
Because $df_q=0$,
$\Hess(f)_q(v,v)=X(Xf)(q)=(f\circ\gamma)''(0)$, which is non-positive
by the one-variable second-derivative test.
\end{proof}

\begin{lemma}[Local smooth extension]
\label{lem:smooth-local-extension}
\lean{MorganTianLib.exists_contMDiff_eventuallyEq_of_contMDiffOn}


If $f$ is smooth on an open $U\subset M$ and $q\in U$, then some smooth
$f'\colon M\to\Ar$ agrees with $f$ near $q$.
\end{lemma}
\begin{proof}
\sketch
\leanok
Multiply $f$ by a bump function supported in $U$ and equal to $1$ near
$q$, then extend by zero.
\end{proof}

\begin{lemma}[Laplacian is non-positive at an interior maximum]
\label{lem:laplacian-nonpositive-at-max}
\lean{MorganTianLib.laplacianAt_nonpos_of_isLocalMaxOn}
\uses{lem:laplacian-local-formula}
\leanok



If a smooth $v$ on an open $U\subset M$ has a local maximum at $q$, then
$dv(q)=0$ and $\Delta v(q)\le0$.
\end{lemma}
\begin{proof}
\sketch
\leanok
\uses{lem:smooth-local-extension,lem:differential-vanishes-at-max,lem:hessian-nonpositive-at-max,lem:laplacian-local-formula}
Extend $v$ smoothly near $q$. Its differential vanishes, and tracing the
negative semidefinite Hessian in an orthonormal basis gives
$\Delta v(q)\le0$.
\end{proof}

\begin{lemma}[Strong maximum principle for smooth subharmonic functions]
\label{lem:hopf-strong-maximum}
\lean{MorganTianLib.hopf_strong_maximum}
\uses{lem:laplacian-christoffel-formula}



Let $U$ be connected and open and let $h\in C^\infty(U)$ satisfy
$\Delta h\ge0$. If $h$ attains $\sup_Uh$, then $h$ is constant.
\end{lemma}
\begin{proof}
\sketch
\uses{lem:laplacian-christoffel-formula,lem:laplacian-nonpositive-at-max}
\leanok
The maximum set is nonempty and closed. To show it is open, work in a
coordinate ball and suppose a maximal point is adjacent to a point where
$h<m$. A Euclidean ball tangent to the maximum set supports the barrier
$$w_\alpha(x)=e^{-\alpha|x-x_0|^2}-e^{-\alpha R^2}.$$
Uniform ellipticity and bounded first-order coefficients give
$\Delta w_\alpha>0$ on an outer annulus when $\alpha$ is large. The
function $h+\delta w_\alpha$ then has neither an interior maximum
(because its Laplacian is positive) nor a maximum on the inner boundary,
while its outward radial derivative at a touching maximal point is
negative. This contradiction makes the maximum set open; connectedness
finishes the proof.
\end{proof}

\begin{theorem}[Maximum principle]
\label{thm:maximum-principle}
\source{morgan-tian}
\uses{rem:laplacian-weak-sense}


Let $f$ be continuous on a connected Riemannian manifold and satisfy
$\Delta f\ge0$ weakly, meaning
$\int f\Delta\varphi\,d\vol\ge0$ for every non-negative compactly
supported smooth $\varphi$. Then $f$ is locally constant near every
local maximum. If it attains a global maximum, it is constant.
\end{theorem}
\begin{proof}
\sketch

\uses{thm:weak-harmonic-regularity,claim:dirichlet-problem-small-balls,claim:weak-subsolution-comparison,rem:laplacian-weak-sense,lem:laplacian-local-formula}
\uses{claim:dirichlet-problem-small-balls,claim:weak-subsolution-comparison,lem:green-identity-compact-support,lem:hopf-strong-maximum,def:exponential-map}
Around a local maximum $p$, choose a small normal ball contained in the
region where $f\le f(p)$. On each smaller sphere solve the Dirichlet
problem with boundary value $f$, obtaining $h$. Green's identity and the
weak comparison principle give $f\le h$ and
$\max h\le f(p)$, so $h(p)=f(p)$. The smooth strong maximum principle
makes $h$ constant, hence $f=f(p)$ on every such sphere and throughout
the ball. The global statement follows because the set of global
maximizers is nonempty, open, and closed.
\end{proof}

\begin{theorem}[Regularity of weakly harmonic functions]
\label{thm:weak-harmonic-regularity}
\uses{rem:laplacian-weak-sense}
\label{eqnlocmax}

If a continuous $f\colon M\to\Ar$ satisfies
$$\int_M f\Delta\varphi\,d\vol=0$$
for every compactly supported smooth $\varphi$, then $f$ is smooth and
classically harmonic.
\end{theorem}
\begin{proof}
\sketch
\uses{claim:dirichlet-problem-small-balls,claim:weak-subsolution-comparison,lem:green-identity-compact-support}
On a sufficiently small normal ball, let $u$ be the harmonic function
with boundary values $f$. Green's identity shows that both $f-u$ and
$u-f$ are weak subsolutions vanishing at the boundary. Comparison gives
$f=u$ on the ball, proving local smoothness and harmonicity.
\end{proof}

\begin{lemma}[First derivatives of the square norm of a gradient]
\label{lem:gradient-norm-sq-gradient}
\lean{MorganTianLib.metricNormSq,MorganTianLib.contMDiff_metricNormSq,MorganTianLib.dir_metricNormSq,MorganTianLib.dir_metricNormSq_gradientField,MorganTianLib.gradientField_metricNormSq_gradientField,MorganTianLib.hessian_metricNormSq_gradientField}
\uses{thm:levi-civita-connection, lem:hessian-symmetric}
\leanok



Let $G=(\nabla f)^*$ for smooth $f$. Then $|\nabla f|^2$ is smooth and
for vector fields $X,Y$,
\begin{enumerate}
\item $Y(|\nabla f|^2)=2\langle\nabla_YG,G\rangle$;
\item $(\nabla|\nabla f|^2)^*=2\nabla_GG$;
\item $\Hess(|\nabla f|^2)(X,Y)=2\langle\nabla_X\nabla_GG,Y\rangle$.
\end{enumerate}
\end{lemma}
\begin{proof}
\sketch
\uses{thm:levi-civita-connection,lem:hessian-symmetric}
\leanok
Metric compatibility proves the first identity. Symmetry of the Hessian
rewrites $\langle\nabla_YG,G\rangle$ as
$\langle\nabla_GG,Y\rangle$, giving the gradient identity. Apply the
gradient formula for the Hessian once more for the third.
\end{proof}

\begin{definition}[Square norm of the Hessian]
\label{def:hessian-norm-sq}
\lean{MorganTianLib.hessianNormSqAt}
\uses{lem:hessian-symmetric}
\leanok



For smooth $f$ and an orthonormal basis $(e_i)$ of $T_pM$, define
$$|\Hess(f)|^2(p)=\sum_{i,j}\Hess(f)_p(e_i,e_j)^2
=g^{ik}g^{jl}\Hess(f)_{ij}\Hess(f)_{kl}.$$
The next lemma shows independence of the basis.
\end{definition}

\begin{lemma}[Basic properties of the Hessian square norm]
\label{lem:hessian-norm-sq-basic}
\lean{MorganTianLib.hessianNormSqAt_eq_sum,MorganTianLib.hessianNormSqAt_nonneg,MorganTianLib.hessianNormSqAt_eq_zero_iff,MorganTianLib.sq_laplacianAt_le_finrank_mul_hessianNormSqAt}
\uses{def:hessian-norm-sq, lem:hessian-symmetric}
\leanok



For $n=\dim M$:
\begin{enumerate}
\item $|\Hess(f)|^2(p)$ is independent of the orthonormal basis;
\item it is non-negative and vanishes exactly when $\Hess(f)_p=0$;
\item $(\Delta f(p))^2\le n|\Hess(f)|^2(p)$.
\end{enumerate}
\end{lemma}
\begin{proof}
\sketch
\uses{lem:hessian-symmetric}
\leanok
Orthogonal change of basis and Parseval preserve the sum of squared
matrix coefficients. The second assertion is the vanishing criterion
for a sum of squares. The third follows from Cauchy--Schwarz applied to
$\Delta f=\sum_i\Hess(f)(e_i,e_i)$.
\end{proof}

\begin{lemma}[Hessian square norm as a trace along the gradient]
\label{lem:hessian-riesz-trace}
\lean{MorganTianLib.metricInner_cov_gradientField_eq_hessianAt,MorganTianLib.sum_metricInner_cov_cov_gradientField_eq_hessianNormSqAt}
\uses{def:hessian-norm-sq, lem:hessian-symmetric}



Let $G=(\nabla f)^*$, let $(e_i)$ be an orthonormal basis at $p$, and
extend it to fields $E_i$. For every field $X$ and $w\in T_pM$,
$$\langle(\nabla_XG)(p),w\rangle=\Hess(f)_p(X(p),w),$$
and
$$\sum_i\langle\nabla_{\nabla_{E_i}G}G,E_i\rangle(p)
=|\Hess(f)|^2(p).$$
\end{lemma}
\begin{proof}
\sketch
\uses{lem:hessian-symmetric}
\leanok
The first identity is the Riesz form of the Hessian-gradient formula.
Writing $\nabla_{E_i}G=\sum_j\Hess(f)(e_i,e_j)e_j$ at $p$ and using
Hessian symmetry turns the trace into the defining sum of squares.
\end{proof}

\begin{lemma}[Smooth local orthonormal frames]
\label{lem:local-orthonormal-frame}
\lean{MorganTianLib.chartFrameNorm_section_contMDiffOn,MorganTianLib.exists_orthonormalFrame,MorganTianLib.orthonormalBasisOfMetricInner}
\uses{def:riemannian-metric}
\leanok



Every $p$ has a neighborhood carrying a smooth orthonormal frame,
whose members may be extended to smooth vector fields on $M$.
\end{lemma}
\begin{proof}
\sketch
\uses{def:riemannian-metric}
\leanok
Apply fibrewise Gram--Schmidt to a coordinate frame. Smoothness follows
because every normalization denominator is positive; multiply by a
bump function equal to $1$ near $p$ to extend the fields globally.
\end{proof}

\begin{lemma}[Derivative of the Laplacian as a trace of the second covariant derivative]
\label{lem:laplacian-trace-commutation}
\lean{MorganTianLib.laplacianAt_eventuallyEq_sum_frame,MorganTianLib.contMDiff_laplacianAt,MorganTianLib.metricInner_secondCov_middle_congr,MorganTianLib.sum_metricInner_secondCov_frame_eq_dir_laplacianAt,MorganTianLib.sum_metricInner_secondCov_gradientField_eq_dir_laplacianAt}
\uses{def:second-covariant-derivative-vector-field, lem:hessian-symmetric}
\leanok



For smooth $f$, $G=(\nabla f)^*$, an orthonormal basis $(e_i)$ at $p$,
and smooth extensions $E_i$, the function $\Delta f$ is smooth and
$$\sum_i\langle\nabla^2G(G,E_i),E_i\rangle(p)=G(\Delta f)(p).$$
\end{lemma}
\begin{proof}
\sketch
\uses{lem:local-orthonormal-frame,lem:second-covariant-derivative-tensorial,lem:hessian-symmetric,lem:hessian-riesz-trace,thm:levi-civita-connection}
\leanok
Use a smooth local orthonormal frame $F_i$, for which
$\Delta f=\sum_i\langle\nabla_{F_i}G,F_i\rangle$. Differentiate along
$G$ and compare with the definition of $\nabla^2G(G,F_i)$. The extra
terms contract the antisymmetric matrix
$\langle\nabla_GF_i,F_j\rangle$ with the symmetric Hessian matrix, so
they vanish.
\end{proof}

\begin{lemma}[Laplacian of the square norm of a one-form]
\label{lem:laplacian-square-norm-one-form}
\uses{thm:levi-civita-connection}

For a smooth one-form $\omega$ and its connection Laplacian,
$$\frac12\Delta|\omega|^2
=\langle\Delta\omega,\omega\rangle+|\nabla\omega|^2.$$
\end{lemma}
\begin{proof}
\sketch
\uses{thm:levi-civita-connection,lem:hessian-symmetric,def:exponential-map,lem:normal-coordinates-expansion}
In normal coordinates at $p$, the metric has zero first derivatives and
the Christoffel symbols vanish. Twice differentiating
$|\omega|^2$ using metric compatibility gives
$$\partial_i\partial_j|\omega|^2
=2\langle\nabla_i\nabla_j\omega,\omega\rangle
+2\langle\nabla_j\omega,\nabla_i\omega\rangle.$$
Trace at $p$ to obtain the identity; tensoriality makes it global.
\end{proof}

\begin{lemma}[Bochner formula for functions]
\label{lem:function-bochner-formula}
\lean{MorganTianLib.laplacianAt_metricNormSq_gradientField,MorganTianLib.function_bochner_formula}
\uses{def:ricci-curvature}
\leanok



For smooth $f$,
$$\frac12\Delta|\nabla f|^2
=|\Hess(f)|^2
+\langle(\nabla f)^*,(\nabla\Delta f)^*\rangle
+\Ric((\nabla f)^*,(\nabla f)^*).$$
\end{lemma}
\begin{proof}
\sketch
\uses{lem:hessian-symmetric,lem:gradient-norm-sq-gradient,def:second-covariant-derivative-vector-field,lem:ricci-commutation-vector-field,lem:ricci-trace-curvature-operator,lem:hessian-riesz-trace,lem:laplacian-trace-commutation,def:hessian-norm-sq,def:ricci-curvature}
\leanok
Set $G=(\nabla f)^*$ and trace
$\Hess(|\nabla f|^2)(X,Y)=2\langle\nabla_X\nabla_GG,Y\rangle$.
Decompose
$$\nabla_{E_i}\nabla_GG
=\nabla^2G(G,E_i)+\mathcal R(E_i,G)G
+\nabla_{\nabla_{E_i}G}G.$$
The three traces are respectively
$G(\Delta f)$, $\Ric(G,G)$, and $|\Hess(f)|^2$ by the preceding
lemmas.
\end{proof}

\begin{corollary}[Bochner vanishing]
\label{cor:bochner-vanishing}
\lean{MorganTianLib.hessianNormSqAt_eq_zero_of_bochner,MorganTianLib.hessianAt_eq_zero_of_bochner,MorganTianLib.cov_gradientField_apply_eq_zero_of_bochner}
\uses{lem:function-bochner-formula, def:hessian-norm-sq, def:ricci-curvature}
\leanok



Suppose $\Delta f$ and $|\nabla f|^2$ are constant and
$\Ric((\nabla f)^*,(\nabla f)^*)(p)\ge0$. Then
$\Hess(f)_p=0$. Hence a harmonic function of constant gradient norm on
a manifold with non-negative Ricci curvature has parallel gradient.
\end{corollary}
\begin{proof}
\sketch
\uses{lem:function-bochner-formula,def:hessian-norm-sq}
\leanok
The derivative terms in Bochner's formula vanish, leaving a sum of the
two non-negative quantities $|\Hess(f)|^2(p)$ and
$\Ric((\nabla f)^*,(\nabla f)^*)(p)$. Thus the Hessian vanishes, and
$\Hess(f)(X,Y)=\langle\nabla_X(\nabla f)^*,Y\rangle$ makes the gradient
parallel.
\end{proof}

\begin{lemma}[Covariant differentiation along smooth maps]
\label{lem:covariant-derivative-along-maps}
\lean{MorganTianLib.chartLocalMap, MorganTianLib.chartFieldCoord_curveVelocity_snd_eventually, MorganTianLib.chartFieldCoord_curveVelocity_fst_eventually, MorganTianLib.curveVelocity_comp_fst_eq, MorganTianLib.curveVelocity_comp_snd_eq, MorganTianLib.hasCovDerivAlongAt_fst_curveVelocity_snd, MorganTianLib.hasCovDerivAlongAt_snd_curveVelocity_fst, MorganTianLib.hasCovDerivAlongAt_fst_snd_symm, MorganTianLib.HasCovDerivAlongAt, MorganTianLib.HasCovDerivAlongAt.unique, MorganTianLib.HasCovDerivAlongAt.add, MorganTianLib.HasCovDerivAlongAt.const_smul, MorganTianLib.HasCovDerivAlongAt.smul_fun, MorganTianLib.hasCovDerivAlongAt_comp, MorganTianLib.HasCovDerivAlongAt.hasDerivAt_metricInner, MorganTianLib.hasGeodesicEquationAt_iff_hasCovDerivAlongAt_velocity_zero, MorganTianLib.IsParallelAlong, MorganTianLib.IsParallelAlong.metricInner_eq, MorganTianLib.IsParallelAlong.apply_eq}
\uses{thm:levi-civita-connection, def:geodesic}
\leanok



Let $\alpha\colon U\subset\Ar^m\to M$ be smooth and let $W$ be a
smooth field along $\alpha$. In local coordinates define
$$\frac{DW}{\partial u^a}
=\left(\partial_aW^k+\Gamma^k_{ij}(\alpha),
\partial_a\alpha^iW^j\right)\partial_k.$$
These expressions agree on overlaps and satisfy:
\begin{enumerate}
\item linearity, the Leibniz rule, and
$D(Z\circ\alpha)/\partial u^a=\nabla_{\partial_a\alpha}Z$;
\item
$\partial_a\langle W_1,W_2\rangle
=\langle D_aW_1,W_2\rangle+\langle W_1,D_aW_2\rangle$;
\item $D_a(\partial_b\alpha)=D_b(\partial_a\alpha)$;
\item a curve is geodesic exactly when $D\gamma'/dt=0$;
\item inner products of parallel fields along a curve are constant, and
parallel fields agreeing once agree everywhere.
\end{enumerate}
\end{lemma}
\begin{proof}
\sketch
\leanok
\uses{thm:levi-civita-connection,def:geodesic}
The coordinate formula is intrinsic because it agrees with
$\nabla_{\partial_a\alpha}Z$ on fields pulled back from coordinate
frames; the Leibniz rule then patches arbitrary fields. Metric
compatibility proves item~2, while equality of mixed partials and
torsion-freeness prove item~3. Item~4 is the coordinate geodesic
equation. Item~5 follows from item~2 and uniqueness for the resulting
linear ODE.
\end{proof}

\begin{lemma}[Covariant differentiation along a curve]
\label{lem:cov-deriv-along-curve}
\lean{MorganTianLib.HasCovDerivAlongAt, MorganTianLib.chartFieldCoord, MorganTianLib.HasCovDerivAlongAt.unique, MorganTianLib.HasCovDerivAlongAt.add, MorganTianLib.HasCovDerivAlongAt.const_smul, MorganTianLib.HasCovDerivAlongAt.sub, MorganTianLib.HasCovDerivAlongAt.smul_fun, MorganTianLib.metricInner_eq_chartMetricInner, MorganTianLib.HasCovDerivAlongAt.hasDerivAt_metricInner, MorganTianLib.IsParallelAlong, MorganTianLib.IsParallelAlong.metricInner_eq, MorganTianLib.IsParallelAlong.apply_eq, MorganTianLib.curveVelocity, MorganTianLib.hasGeodesicEquationAt_iff_hasCovDerivAlongAt_velocity_zero, MorganTianLib.fieldChartRep, MorganTianLib.hasCovDerivAlongAt_comp, MorganTianLib.cov_apply_eq_fderiv_add_chartChristoffelContraction, MorganTianLib.hasCovDerivAlongAt_comp_cov, MorganTianLib.hasCovDerivAlongAt_comp_zero, MorganTianLib.isParallelAlong_comp_of_cov_eq_zero, MorganTianLib.isParallelAlong_curveVelocity_of_isGeodesic}
\uses{thm:levi-civita-connection, def:geodesic}
\leanok



For a field $W$ along a curve $\gamma$ and a chart state $(u,V)$ near
$t_0$, define
$$\frac{DW}{dt}(t_0)=\dot V(t_0)
+\Gamma(\dot u(t_0),W(t_0))(u(t_0)).$$
This is unique, additive, homogeneous, and obeys the scalar Leibniz
rule. For a smooth field $Z$ on $M$ it equals
$\nabla_{\gamma'(t_0)}Z$. It is compatible with the metric, the geodesic
equation is $D\gamma'/dt=0$, inner products of parallel fields are
constant, and two parallel fields agreeing once agree everywhere.
\end{lemma}
\begin{proof}
\sketch
\leanok
\uses{thm:levi-civita-connection,def:geodesic,lem:covariant-derivative-along-maps}
The chain and product rules applied to the chart formula prove the
algebraic identities. Expanding a smooth field in a chart frame gives
the standard coordinate formula for $\nabla_XZ$, proving the intrinsic
interpretation. Metric compatibility gives the derivative of inner
products, and the formula applied to $W=\gamma'$ is precisely the
geodesic equation. Constancy and uniqueness of parallel fields follow
from the metric identity and positive definiteness.
\end{proof}

\begin{lemma}[Flow of a parallel gradient field]
\label{lem:parallel-gradient-flow}
\lean{MorganTianLib.isParallelAlong_gradientField_comp_of_bochner, MorganTianLib.curveVelocity_eq_gradientField_of_bochner, MorganTianLib.metricInner_curveVelocity_self_of_bochner, MorganTianLib.hasDerivAt_comp_of_bochner, MorganTianLib.comp_eq_add_mul_of_bochner, MorganTianLib.isMIntegralCurve_gradientField_of_bochner, MorganTianLib.isMIntegralCurve_smoothVectorField_eq, MorganTianLib.isMIntegralCurve_smoothVectorField_comp_add, MorganTianLib.IsContGeodesicallyComplete, MorganTianLib.isContGeodesicallyComplete_of_complete, MorganTianLib.exists_isMIntegralCurve_gradientField_of_complete, MorganTianLib.exists_gradientFlowLine_of_bochner, MorganTianLib.smoothVectorFieldFlow, MorganTianLib.smoothVectorFieldFlow_zero, MorganTianLib.smoothVectorFieldFlow_add, MorganTianLib.smoothVectorFieldFlow_bijective, MorganTianLib.continuous_smoothVectorFieldFlow, MorganTianLib.smoothVectorFieldFlowHomeomorph, MorganTianLib.comp_smoothVectorFieldFlow_gradientField_of_bochner, MorganTianLib.fderiv_fieldChartRep_gradientField_of_bochner, MorganTianLib.chartMetricInner_variational_eq_left, MorganTianLib.exists_localFlow_hasStrictFDerivAt, MorganTianLib.metricPreserving_of_localFlowRep, MorganTianLib.exists_flowIsometryBoxAt, MorganTianLib.metricPreserving_smoothVectorFieldFlow_of_bochner, HasStrictFDerivAt.continuousAt_derivMap, MorganTianLib.contDiffOn_one_of_hasStrictFDerivAt, MorganTianLib.contMDiffAt_smoothVectorFieldFlow_of_bochner, MorganTianLib.contMDiff_smoothVectorFieldFlow_of_bochner, MorganTianLib.exists_localFlow_hasStrictFDerivAt_uncurry, MorganTianLib.exists_flowJointC1BoxAt, MorganTianLib.contMDiffAt_smoothVectorFieldFlow_uncurry_of_bochner, MorganTianLib.contMDiff_smoothVectorFieldFlow_uncurry_of_bochner, MorganTianLib.mfderiv_smoothVectorFieldFlow_gradientField_of_bochner, MorganTianLib.pathELength_comp_eq_of_metricPreserving, MorganTianLib.edist_smoothVectorFieldFlow_le_of_bochner, MorganTianLib.isometry_smoothVectorFieldFlow_of_bochner}
\uses{thm:levi-civita-connection, def:geodesic, thm:hopf-rinow, def:exponential-map}


Let $(X,g)$ be complete and let $B\colon X\to\Ar$ be smooth with
$\Hess(B)=0$ and $|\nabla B|=1$. Put $V=(\nabla B)^*$. Then:
\begin{enumerate}
\item $V$ is parallel;
\item $V$ is complete and its smooth flow is
$$\theta_t(x)=\exp_x(tV(x)),$$
whose trajectories are unit-speed geodesics, with
$\theta_0=\mathrm{id}$ and $\theta_s\theta_t=\theta_{s+t}$;
\item $B(\theta_t(x))=B(x)+t$;
\item every $\theta_t$ is an isometry.
\end{enumerate}
\end{lemma}
\begin{proof}
\sketch
\uses{thm:levi-civita-connection,lem:hessian-symmetric,def:geodesic,thm:hopf-rinow,def:exponential-map,lem:covariant-derivative-along-maps}
The identity
$\langle\nabla_WV,Y\rangle=\Hess(B)(W,Y)$ proves parallelism.
Completeness of $X$ extends the geodesic with initial velocity $V(x)$
for all time; its velocity and $V$ are parallel and initially equal, so
the geodesic is the integral curve of $V$. ODE uniqueness yields the
flow law and inverse $\theta_{-t}$. Since $dB(V)=1$, integration gives
item~3. For a flow variation $\alpha(s,t)=\theta_t(\sigma(s))$, mixed
covariant derivatives show $D_t\partial_s\alpha=0$; hence
$|d\theta_t(w)|=|w|$, proving item~4 by polarization.
\end{proof}

\begin{lemma}[Level sets of a function with parallel unit gradient]
\label{lem:parallel-gradient-level-sets}
\lean{MorganTianLib.nonempty_preimage_of_bochner, MorganTianLib.isClosed_preimage_of_continuous, MorganTianLib.metricInner_curveVelocity_gradientField_of_bochner, MorganTianLib.hasDerivAt_comp_zero_of_bochner, MorganTianLib.comp_eq_of_bochner_of_metricInner_eq_zero, MorganTianLib.mapsTo_preimage_of_bochner_of_metricInner_eq_zero, MorganTianLib.metricInner_gradientField_eq_zero_of_mfderiv_eq_zero, MorganTianLib.isConnected_preimage_of_bochner, MorganTianLib.levelSetFlowIsometry, MorganTianLib.exists_openPartialHomeomorph_comp_symm_eq_affine, MorganTianLib.hasFDerivAt_comp_extChartAt_symm, MorganTianLib.mfderiv_ne_zero_of_metricNormSq_gradientField_eq_one, MorganTianLib.exists_extChartAt_openPartialHomeomorph_comp_symm_eq_affine, MorganTianLib.levelDifferential, MorganTianLib.levelRieszVector, MorganTianLib.mem_orthogonal_levelRieszVector_iff, MorganTianLib.levelHyperplaneBasis, MorganTianLib.sliceProj, MorganTianLib.sliceEmb, MorganTianLib.adaptedStraightening, MorganTianLib.levelSliceChart, MorganTianLib.levelSetChartedSpace, MorganTianLib.contDiffOn_levelSliceChart_trans, MorganTianLib.isManifold_levelSet, MorganTianLib.contMDiff_levelSet_val, MorganTianLib.mfderiv_levelSet_val_injective, MorganTianLib.range_mfderiv_levelSet_val, MorganTianLib.exists_hasCovDerivAlongAt_curveVelocity, MorganTianLib.metricInner_gradientField_curveVelocity_eq_zero_of_level, MorganTianLib.metricInner_gradientField_covDerivVelocity_eq_zero_of_level, MorganTianLib.levelSetInducedMetric, MorganTianLib.levelSetInducedMetric_metricInner, MorganTianLib.metricInner_mfderiv_levelSet_val, MorganTianLib.exists_hasCovDerivAlongAt_curveVelocity_levelSet}
\uses{lem:parallel-gradient-flow, def:geodesic, def:riemannian-metric}


Under the hypotheses of Lemma~\ref{lem:parallel-gradient-flow}, let
$N_c=B^{-1}(c)$ with the induced metric. Then:
\begin{enumerate}
\item $N_c$ is a nonempty closed embedded hypersurface,
$T_yN_c=\ker dB_y=V(y)^\perp$, and $V$ is its unit normal. Locally
$B$ is the final coordinate and $N_c$ the corresponding slice; every
smooth map into $X$ with image in $N_c$ is smooth as a map into $N_c$;
\item an $X$-geodesic initially tangent to $N_c$ remains in $N_c$;
\item ambient and intrinsic covariant acceleration agree for curves in
$N_c$. Thus $N_c$ is totally geodesic.
\end{enumerate}
\end{lemma}
\begin{proof}
\sketch
\uses{lem:parallel-gradient-flow,lem:covariant-derivative-along-maps,def:geodesic,def:exponential-map,thm:levi-civita-connection}
The flow moves every point to level $c$, proving nonemptiness. Since
$dB(V)=1$, the regular-level theorem supplies slice charts and
$T_yN_c=V(y)^\perp$. Along a geodesic,
$\langle V,\gamma'\rangle$ is constant because both fields are parallel;
zero initial value proves item~2. For a curve in $N_c$, differentiating
$\langle V,\gamma'\rangle=0$ makes its ambient acceleration tangent.
The first-variation identity
$$E'(0)=-\int\left\langle Y,\frac{D\gamma'}{dt}\right\rangle dt$$
has the same left side for the ambient and induced metrics. Testing with
compactly supported tangent variation fields shows ambient and intrinsic
accelerations coincide, proving item~3.
\end{proof}

\begin{proposition}[Splitting along a parallel unit gradient]
\label{prop:parallel-gradient-splitting}
\lean{MorganTianLib.bochnerSplittingHomeomorph, MorganTianLib.bochnerSplittingHomeomorphOfComplete, MorganTianLib.mfderivReal, MorganTianLib.abs_metricInner_le_sqrt_mul_sqrt, MorganTianLib.ofReal_abs_sub_le_pathELength, MorganTianLib.ofReal_abs_sub_le_edist_of_bochner, MorganTianLib.abs_sub_le_dist_of_bochner, MorganTianLib.edist_smoothVectorFieldFlow_self_of_bochner, MorganTianLib.edist_smoothVectorFieldFlow_flow_of_bochner, MorganTianLib.isMinGeodesicOn_smoothVectorFieldFlow_of_bochner, MorganTianLib.edist_smoothVectorFieldFlow_pair_le_of_bochner, MorganTianLib.ofReal_abs_le_edist_smoothVectorFieldFlow_pair_of_bochner, MorganTianLib.edist_le_edist_smoothVectorFieldFlow_pair_add_of_bochner, MorganTianLib.isComplete_preimage_of_continuous, MorganTianLib.isGeodesicRay_smoothVectorFieldFlow_of_bochner, MorganTianLib.busemann_smoothVectorFieldFlow_of_bochner, MorganTianLib.sub_le_busemann_smoothVectorFieldFlow_of_bochner, MorganTianLib.mfderiv_smoothVectorFieldFlow_comp_apply_one_of_bochner, MorganTianLib.enorm_mfderiv_smoothVectorFieldFlow_comp_of_bochner, MorganTianLib.pathELength_smoothVectorFieldFlow_comp_of_bochner, MorganTianLib.pathELength_smoothVectorFieldFlow_comp_levelSet_of_bochner, MorganTianLib.hasDerivAt_comp_gradientField, MorganTianLib.enorm_mfderiv_eq_sqrt_metricInner_of_hasMFDerivAt, MorganTianLib.enorm_mfderiv_levelProjection_of_bochner, MorganTianLib.edist_levelProjection_le_of_bochner, MorganTianLib.ofReal_mul_add_mul_le_edist_smoothVectorFieldFlow_pair_of_bochner, MorganTianLib.ofReal_sqrt_le_edist_smoothVectorFieldFlow_pair_of_bochner, MorganTianLib.continuous_metricInner_mfderiv_of_contMDiff, MorganTianLib.pathELength_levelProjection_le_of_bochner, MorganTianLib.pathELength_eq_ofReal_integral_of_contMDiff, MorganTianLib.edist_smoothVectorFieldFlow_levelPath_le_of_bochner, MorganTianLib.edist_smoothVectorFieldFlow_pair_le_ofReal_sqrt_of_bochner, MorganTianLib.edist_smoothVectorFieldFlow_pair_eq_of_bochner, MorganTianLib.dist_smoothVectorFieldFlow_pair_of_bochner, MorganTianLib.busemann_smoothVectorFieldFlow_eq_of_bochner, MorganTianLib.prod_dist_eq_sqrt_of_L2, MorganTianLib.dist_eq_sqrt_dist_levelProjection_of_bochner, MorganTianLib.bochnerSplittingIsometry, MorganTianLib.metricInner_flowVariation_of_bochner, MorganTianLib.levelSetInducedMetric, MorganTianLib.levelSetProductMetric, MorganTianLib.DCProductMetric_metricInner_mk, MorganTianLib.levelSetProductMetric_metricInner, MorganTianLib.levelSetProductMetric_metricInner_ambient, MorganTianLib.mfderiv_prod_eq_add, MorganTianLib.mfderiv_bochnerSplittingMap_apply, MorganTianLib.metricInner_mfderiv_bochnerSplittingMap, MorganTianLib.mfderiv_bochnerSplittingMap_injective}
\uses{def:geodesic, thm:hopf-rinow, def:exponential-map, thm:levi-civita-connection, def:riemannian-metric}


Let $(X,g)$ be complete and connected and let $B\colon X\to\Ar$ be
smooth with $\Hess(B)=0$ and $|\nabla B|=1$. Then
$N=B^{-1}(0)$ is a nonempty connected complete totally geodesic embedded
hypersurface, and
$$\Phi\colon(N\times\Ar,g_N\oplus dt^2)\to(X,g),qquad
\Phi(y,t)=\exp_y(t(\nabla B)^*(y)),$$
is an isometry onto $X$, where
$(g_N\oplus dt^2)((w,a),(w',a'))=g_N(w,w')+aa'$.
\end{proposition}
\begin{proof}
\sketch
\uses{lem:parallel-gradient-flow,lem:parallel-gradient-level-sets,thm:levi-civita-connection,lem:hessian-symmetric,def:riemannian-metric}
Let $\theta$ be the gradient flow. Its inverse-coordinate map is
$$\Psi(x)=(\theta_{-B(x)}(x),B(x)),$$
so $\Phi(y,t)=\theta_t(y)$ is a diffeomorphism with inverse $\Psi$.
Connectedness of $X$ then implies connectedness of $N$. For
$w,w'\in T_yN$, the flow is isometric,
$d\Phi(0,1)=V(\theta_t(y))$, and $dB(w)=0$; consequently
$$\langle d\Phi(w,a),d\Phi(w',a')\rangle=g_N(w,w')+aa'.$$
Thus $\Phi$ is a Riemannian isometry. Finally, the induced distance on
$N$ equals the restriction of $d_X$: projecting any product path onto
$N$ does not increase length. A Cauchy sequence in $N$ therefore
converges in $X$, and closedness of $N$ places the limit in $N$.
\end{proof}

\begin{lemma}[Busemann functions along a minimizing line]
\label{lem:busemann-along-line}
\lean{MorganTianLib.busemann_apply_line_pair}
\source{morgan-tian}
\uses{def:busemann-function, def:minimizing-geodesic-window}
\leanok




For a minimizing line $\lambda\colon\Ar\to M$, let $B_+$ and $B_-$
be the Busemann functions of $t\mapsto\lambda(t)$ and
$t\mapsto\lambda(-t)$. Then
$$B_+(\lambda(u))=-u,qquad B_-(\lambda(u))=u\qquad(u\in\Ar).$$
\end{lemma}
\begin{proof}
\sketch
\leanok
\uses{lem:busemann-le-approx}
For $t\ge\max(u,0)$,
$B_{\lambda_+,t}(\lambda(u))=(t-u)-t=-u$, while for all $t\ge0$ the
same approximant is $|t-u|-t\ge-u$. Apply this also to the reversed
line.
\end{proof}

\begin{lemma}[The Busemann pair of a line is non-negative]
\label{lem:busemann-pair-nonneg}
\lean{MorganTianLib.busemann_add_busemann_neg_nonneg}
\source{morgan-tian}
\uses{def:busemann-function, def:minimizing-geodesic-window}
\leanok




For the Busemann pair of a minimizing line,
$$B_+(x)+B_-(x)\ge0\qquad(x\in M).$$
\end{lemma}
\begin{proof}
\sketch
\leanok
\uses{def:busemann-approx}
For $s,t\ge0$, the triangle inequality between $\lambda(-s)$,
$x$, and $\lambda(t)$ gives
$B_{\lambda_+,t}(x)+B_{\lambda_-,s}(x)\ge0$. Take both infima.
\end{proof}

\begin{corollary}[The Busemann pair vanishes on its line]
\label{cor:busemann-pair-vanishes-on-line}
\lean{MorganTianLib.busemann_add_busemann_neg_apply_line}
\uses{lem:busemann-along-line}
\leanok



$B_+(\lambda(u))+B_-(\lambda(u))=0$ for every $u\in\Ar$.
\end{corollary}
\begin{proof}
\sketch
\leanok
\uses{lem:busemann-along-line}
Use the two identities of Lemma~\ref{lem:busemann-along-line}.
\end{proof}

\begin{lemma}[Minimizing line implies isometric product]
\label{lem:minimizing-line-implies-product}
\source{morgan-tian}
\uses{def:ricci-curvature}
\label{line}


A connected complete Riemannian manifold of non-negative Ricci
curvature that contains a minimizing line is isometric to
$N\times\Ar$ for some Riemannian manifold $N$
\cite{CheegerGromoll2}.
\end{lemma}
\begin{proof}
\sketch
\uses{def:busemann-function,rem:laplacian-weak-sense,prop:busemann-laplacian,thm:maximum-principle,thm:weak-harmonic-regularity,def:asymptotic-ray,lem:asymptotic-ray-exists,lem:asymptotic-ray,lem:busemann-gradient-norm-one,lem:function-bochner-formula,prop:parallel-gradient-splitting,lem:busemann-along-line,lem:busemann-pair-nonneg,cor:busemann-pair-vanishes-on-line}
Let $B_\pm$ be the Busemann functions of the two rays of the line.
Both are weakly superharmonic, while $B_++B_-\ge0$ and vanishes on the
line. The weak maximum principle applied to $-(B_++B_-)$ gives
$B_-=-B_+$. Thus $B_+$ is weakly harmonic, hence smooth and harmonic by
regularity. Through each $x$ there is an asymptotic ray, so
Lemma~\ref{lem:busemann-gradient-norm-one} gives
$|\nabla B_+|=1$. Bochner's formula and $\Ric\ge0$ then imply
$\Hess(B_+)=0$. Proposition~\ref{prop:parallel-gradient-splitting}
applied to $B_+$ yields the product isometry.
\end{proof}

\begin{theorem}[Splitting Theorem]
\label{thm:splitting-theorem}
\source{morgan-tian}
\uses{def:ricci-curvature, def:space-of-ends}
\label{splitting}


If $M$ is complete, has non-negative Ricci curvature, and has at least
two ends, then $M$ is isometric to $N\times\Ar$ for a compact manifold
$N$.
\end{theorem}
\begin{proof}
\sketch
\uses{lem:ends-exist,lem:minimizing-line-implies-product,def:space-of-ends,lem:compact-set-in-codim-zero-submanifold}
The ends hypothesis implies that $M$ is connected and noncompact.
Lemma~\ref{lem:ends-exist} gives a minimizing line, and
Lemma~\ref{lem:minimizing-line-implies-product} gives
$M\cong N\times\Ar$ with $N$ connected. If $N$ were noncompact, then
for every compact $K\subset N\times\Ar$ one could choose
$C\times[-a,a]\supset K$ and $z\in N\setminus C$; paths through
$(z,a+1)$ show that the complement of this larger compact set is
connected. Cofinality of compact codimension-zero submanifolds then
forces all ends of $N\times\Ar$ to agree. This contradicts invariance of
the at-least-two ends of $M$ under the product isometry. Hence $N$ is
compact.
\end{proof}

\section{$\epsilon$-necks}

Fix $0<\epsilon<1/2$ and put $k=\lfloor\epsilon^{-1}\rfloor$.

\begin{definition}[Epsilon-closeness in $C^{[1/\epsilon]}$-topology]
\label{def:epsilon-close-topology}
\lean{MorganTianLib.epsilonDerivativeOrder,MorganTianLib.iteratedCovariantMetricDifference,MorganTianLib.metricCovariantDerivativeNormSq,MorganTianLib.epsilonClosenessDensity,MorganTianLib.EpsilonClose,MorganTianLib.EpsilonCloseFamily}
\uses{thm:levi-civita-connection}
\leanok
\label{epsclose}

Let $g_0$ be a metric on $M$ and $X\subset M$ open. A metric $g$ on
$X$ is \emph{$\epsilon$-close to $g_0$ in $C^{[1/\epsilon]}$} if
\begin{equation}\label{eqnepsclose}
\sup_{x\in X}\left(|g-g_0|_{g_0}^2
+\sum_{\ell=1}^k|\nabla_{g_0}^\ell g|_{g_0}^2\right)<\epsilon^2.
\end{equation}
For families $g(t),g_0(t)$ on an interval $I$, the corresponding
condition is the same supremum over $(x,t)\in X\times I$, with the
metric, covariant derivatives, and norms at time $t$.
\end{definition}

\begin{remark}[Ricci flow and metric topology]
\label{rem:ricci-flow-metric-topology}
\uses{def:epsilon-close-topology}

Closeness of metric families in this sense imposes no time-derivative
bound. For Ricci flows, $C^{2k}$-closeness is equivalent to control of
the $r$th time derivatives of the $s$th covariant derivatives whenever
$s+2r\le2k$.
\end{remark}

\begin{definition}[Epsilon-neck structure]
\label{def:epsilon-neck}
\lean{MorganTianLib.roundSphereMetric, MorganTianLib.IsRoundCylinderMetric, MorganTianLib.EpsilonNeckStructure, MorganTianLib.epsilonNeckSlice, MorganTianLib.EpsilonNeckStructure.axialCoordinate, MorganTianLib.EpsilonNeckStructure.axialVector, MorganTianLib.EpsilonNeckStructure.negativeEnd, MorganTianLib.EpsilonNeckStructure.positiveEnd, MorganTianLib.EpsilonNeckStructure.fractionalSubneck, MorganTianLib.EpsilonNeckStructure.reversedParametrization, MorganTianLib.EpsilonNeckStructure.reversedParametrization_apply, MorganTianLib.AmbientEpsilonNeck, MorganTianLib.AmbientEpsilonNeck.ambientCenter, MorganTianLib.EpsilonNeckStructure.scale, MorganTianLib.AmbientEpsilonNeck.scale}
\uses{def:sectional-curvature, def:epsilon-close-topology}
\leanok
\label{epsneck}

An \emph{$\epsilon$-neck structure} on a Riemannian $3$-manifold
$(N,g)$ centered at $x$ is a diffeomorphism
$$\varphi\colon S^2\times(-\epsilon^{-1},\epsilon^{-1})\to N$$
with $x\in\varphi(S^2\times\{0\})$ such that
$R(x)\varphi^*g$ is $\epsilon$-close in $C^{[1/\epsilon]}$ to the
product of the Euclidean interval and the round $S^2$ of Gaussian
curvature $1/2$. The slices are the neck's $2$-spheres and the slice at
$0$ is central. Define
$$s_N=p_2\circ\varphi^{-1},qquad
\frac{\partial}{\partial s_N}=\varphi_*\frac{\partial}{\partial s}.$$
The signs of $s_N$ specify the two ends and all fractional subnecks; the
reversed structure is obtained from
$\varphi\circ(\mathrm{id}_{S^2}\times(-1))$. An $\epsilon$-neck in an
ambient manifold is such a codimension-zero submanifold. Its scale is
$$r_N=R(x)^{-1/2}.$$
\end{definition}

\begin{lemma}[Epsilon-neck separates soul from end]
\label{lem:epsilon-neck-separates}
\uses{def:epsilon-neck, lem:metric-sphere-meets-neck-topologically}
\label{neckseparate}

Assume $\epsilon$ is small enough for Lemma~\ref{topiso}. Let $M$ be a
noncompact positively curved $3$-manifold with soul $p$. The central
sphere of every $\epsilon$-neck $N$ disjoint from $p$ separates $p$
from the end. If two such necks are mutually disjoint, their central
spheres bound a region diffeomorphic to $S^2\times I$.
\end{lemma}
\begin{proof}
\sketch
\uses{lem:metric-sphere-meets-neck-topologically}
For a point $z$ in the middle third of $N$, Lemma~\ref{topiso}
identifies the metric sphere $\partial B(p,d(p,z))$ with a sphere in
$N$ isotopic to its central sphere. This sphere separates the soul from
infinity. Applying this to two disjoint necks gives the product region.
\end{proof}

\begin{proposition}[Positively curved manifolds don't have small epsilon-necks]
\label{prop:epsilon-neck-narrow}
\uses{def:epsilon-neck, def:sectional-curvature}
\label{narrows}

For all sufficiently small $\epsilon>0$, a complete positively curved
Riemannian $3$-manifold contains no $\epsilon$-necks of arbitrarily
small scale.
\end{proposition}
\begin{proof}
  \uses{thm:soul-theorem, prop:busemann-laplacian, lem:neck-geodesic-directions, def:epsilon-neck, lem:metric-sphere-meets-neck-topologically}
\sketch

The compact case is immediate. Otherwise let $p$ be the point soul and
note that every positive level of $d(p,\cdot)$ is connected. Choose two
disjoint necks $N_1,N_2$ away from it, ordered toward the
end by Lemma~\ref{lem:epsilon-neck-separates}. Construct a non-negative
cutoff $\psi$ equal to $1$ between their central regions and changing
only along their middle thirds. If $B_\lambda$ is the Busemann function
of a ray from $p$, the neck direction estimate gives
$$\int_M\langle\nabla B_\lambda,\nabla\psi\rangle\,d\vol
=\bigl(\alpha_2R(x_2)^{-1}-\alpha_1R(x_1)^{-1}\bigr)
\Vol_{h_0}(S^2),$$
where $\alpha_i\to1$ as $\epsilon\to0$. Weak superharmonicity of
$B_\lambda$ makes the integral non-negative. For sufficiently small
$\epsilon$, this implies $R(x_2)\le2R(x_1)$ for any ordered pair, ruling
out necks whose scales tend to zero.
\end{proof}

\begin{corollary}[Scale bound between separated $\epsilon$-necks]
\label{cor:epsilon-neck-scale-bound}
\uses{def:epsilon-neck, thm:soul-theorem, lem:epsilon-neck-separates, lem:metric-sphere-meets-neck-topologically}
\label{strnarrows}

Fix sufficiently small $\epsilon$. There is $C(\epsilon)<\infty$ such
that if $N$ is an $\epsilon$-neck centered at $x$ in a noncompact
positively curved $3$-manifold, disjoint from a soul $p$, and $N'$ is
an $\epsilon$-neck centered at $x'$ separated from $p$ by $N$, then
$$R(x')\le C(\epsilon)R(x).$$
\end{corollary}

\section{Forward difference quotients}

\begin{definition}[Forward difference quotient bounds]
\label{def:forward-difference-quotient}
\lean{MorganTianLib.ForwardDiffQuotientLE}
\source{morgan-tian}
\leanok



For continuous $f\colon[a,b]\to\Ar$, its forward difference quotient
at $t\in[a,b)$ is at most $c$ if
$$\limsup_{\triangle t\to0^+}
\frac{f(t+\triangle t)-f(t)}{\triangle t}\le c;$$
equivalently, for every $r>c$, all sufficiently small positive
$\triangle t$ make this quotient less than $r$. The lower bound is
defined analogously using $\liminf$.
\end{definition}

\begin{lemma}[Weakening a forward difference quotient bound]
\label{lem:forward-difference-quotient-mono}
\lean{MorganTianLib.ForwardDiffQuotientLE.mono}
\uses{def:forward-difference-quotient}
\leanok



If the forward difference quotient of $f$ at $t$ is at most $c$ and
$c\le c'$, then it is at most $c'$.
\end{lemma}
\begin{proof}
\sketch
\leanok
Every $r>c'$ also satisfies $r>c$.
\end{proof}

\begin{lemma}[Right derivatives bound forward difference quotients]
\label{lem:right-derivative-forward-difference-quotient}
\lean{MorganTianLib.HasDerivWithinAt.forwardDiffQuotientLE}
\uses{def:forward-difference-quotient}
\leanok



If $f$ has right derivative $c$ at $t$, then its forward difference
quotient there is at most $c$.
\end{lemma}
\begin{proof}
\sketch
\leanok
The relevant difference quotients converge to $c$.
\end{proof}

\begin{lemma}[Forward difference quotient comparison lemma]
\label{lem:forward-difference-quotient}
\lean{MorganTianLib.le_of_forwardDiffQuotientLE}
\source{morgan-tian}
\uses{def:forward-difference-quotient}
\leanok
\label{fordiffquot}




Let $f\colon[a,b]\to\Ar$ be continuous and let
$\psi\in C^1([a,b]\times\Ar)$. Suppose the forward difference quotient
of $f$ satisfies
$$\frac{df}{dt}(t)\le\psi(t,f(t))\qquad(t\in[a,b)),$$
and $G'=\psi(t,G)$ on $[a,b]$ with $f(a)\le G(a)$. Then
$f(t)\le G(t)$ for all $t\in[a,b]$.
\end{lemma}
\begin{proof}
\sketch
\uses{def:forward-difference-quotient}
\leanok
On the compact range of $f$ and $G$, choose a Lipschitz constant $K$
for $\psi$ in its second variable. For small $\epsilon>0$, set
$$B_\epsilon(t)=G(t)+\epsilon e^{2K(t-a)}.$$
If $f$ first touched this strict supersolution from below, then at the
touching time
$\psi(t,f(t))<B_\epsilon'(t)$. The forward quotient definition and the
right derivative of $B_\epsilon$ therefore keep $f<B_\epsilon$ just
after that time, a contradiction. Thus $f\le B_\epsilon$; let
$\epsilon\downarrow0$.
\end{proof}

\begin{lemma}[The maximum of a compact family is continuous]
\label{lem:max-function-continuous-product}
\lean{MorganTianLib.continuous_iSup_of_compact}
\leanok


If $X$ is compact and $F\colon X\times\Ar\to\Ar$ is continuous, then
$$F_{\max}(t)=\sup_{x\in X}F(x,t)$$
is continuous.
\end{lemma}
\begin{proof}
\sketch
\leanok
For nonempty $X$, a maximizer at $t_0$ gives lower semicontinuity.
Finitely many product neighborhoods covering $X\times\{t_0\}$ give the
matching upper bound uniformly in $x$. The empty case is immediate.
\end{proof}

\begin{lemma}[Forward difference quotient of the maximum of a compact family]
\label{lem:forward-difference-maximum-product}
\lean{MorganTianLib.forwardDiffQuotientLE_iSup}
\uses{def:forward-difference-quotient, lem:max-function-continuous-product}
\leanok



Let $X$ be nonempty and compact, let $t_0<b$, and let $F,F'$ be
continuous on $X\times\Ar$, with
$\partial_sF(x,s)=F'(x,s)$ for $s\in(t_0,b)$. If
$F'(x,t_0)\le c$ at every maximizer of $F(\cdot,t_0)$, then the
forward difference quotient of $F_{\max}$ at $t_0$ is at most $c$.
\end{lemma}
\begin{proof}
\sketch
\leanok
\uses{lem:max-function-continuous-product}
Otherwise some $r>c$ and times $z\downarrow t_0$ have maximum
difference quotient at least $r$. Choose maximizers $x_z$ at time $z$.
The mean value theorem produces $s_z\in(t_0,z)$ with
$F'(x_z,s_z)\ge r$. Compactness gives a cluster point $x_0$, and
continuity makes $x_0$ a maximizer at $t_0$ with
$F'(x_0,t_0)\ge r>c$, a contradiction.
\end{proof}

\begin{corollary}[Forward difference maximum property, product case]
\label{cor:forward-difference-maximum-product}
\lean{MorganTianLib.iSup_le_of_forwardDiff_max}
\uses{def:forward-difference-quotient}



Let $X$ be nonempty and compact, and let $F,F'$ be continuous on
$X\times\Ar$, with $\partial_sF(x,s)=F'(x,s)$. Suppose
$$F'(x,t)\le\psi(t,F_{\max}(t))$$
whenever $x$ maximizes $F(\cdot,t)$, where
$\psi\in C^1([a,b]\times\Ar)$. If $G'=\psi(t,G)$ and
$F_{\max}(a)\le G(a)$, then $F_{\max}(t)\le G(t)$ on $[a,b]$.
\end{corollary}
\begin{proof}
\sketch
\leanok
\uses{lem:max-function-continuous-product,lem:forward-difference-maximum-product,lem:forward-difference-quotient}
The preceding lemmas give continuity of $F_{\max}$ and its required
forward quotient bound. Apply
Lemma~\ref{lem:forward-difference-quotient}.
\end{proof}

\begin{lemma}[Flow of a vector field near a compact set]
\label{lem:vector-field-flow-near-compact}
\lean{MorganTianLib.exists_forall_hasDerivWithinAt_continuousOn_of_contDiffAt, MorganTianLib.exists_localFlowAt, MorganTianLib.exists_localFlow_of_isCompact, MorganTianLib.exists_isLocalFlow_of_isCompact}
\leanok


Let $\chi$ be a smooth vector field on a boundaryless Hausdorff smooth
manifold and let $Z$ be compact. There are $\eta>0$, an open
$U\supset Z$, and a jointly continuous
$\Phi(x,s)=\varphi_s(x)$ on $U\times(-\eta,\eta)$ such that
$$\varphi_0(x)=x,qquad \partial_s\varphi_s(x)=\chi(\varphi_s(x)).$$
Consequently
$$\partial_sF(\varphi_s(x))=\chi(F)(\varphi_s(x)),$$
and if $\chi({\bf t})=1$, then
${\bf t}(\varphi_s(x))={\bf t}(x)+s$.
\end{lemma}
\begin{proof}
\sketch
\leanok
In a chart, Picard iteration gives a common short-time interval for
initial data in a smaller closed ball. Gr\"onwall yields continuous
dependence on the initial point. Cover $Z$ by finitely many such flow
boxes and take the minimum lifetime; uniqueness glues them on overlaps.
The final identities follow from the chain rule and integration.
\end{proof}

\begin{lemma}[Forward difference maximum property along a local flow]
\label{lem:forward-difference-maximum-flow}
\lean{MorganTianLib.IsLocalFlow, MorganTianLib.continuousOn_levelMax, MorganTianLib.forwardDiffQuotientLE_levelMax, MorganTianLib.levelMax_le_of_isLocalFlow}
\uses{def:forward-difference-quotient}
\leanok



Let ${\bf t},F,\chi F\colon M\to\Ar$, with ${\bf t}$ and $\chi F$
continuous, and let $\mathcal Z$ be the set of fibrewise maximizers of
$F$. Assume:
\begin{enumerate}
\item $\mathcal Z$ is compact and meets every fibre
${\bf t}^{-1}(t)$ for $t\in[a,b]$;
\item a jointly continuous local flow $\Phi$ near $\mathcal Z$ satisfies
$$\Phi(x,0)=x,quad {\bf t}(\Phi(x,s))={\bf t}(x)+s,quad
\partial_sF(\Phi(x,s))=\chi F(\Phi(x,s));$$
\item $\chi F(x)\le\psi({\bf t}(x),F(x))$ on $\mathcal Z$, with
$\psi\in C^1([a,b]\times\Ar)$.
\end{enumerate}
Then $F_{\max}(t)=\max_{{\bf t}(x)=t}F(x)$ is continuous and has
forward difference quotient at most $\psi(t,F_{\max}(t))$. If
$G'=\psi(t,G)$ and $F_{\max}(a)\le G(a)$, then
$F_{\max}\le G$ on $[a,b]$.
\end{lemma}
\begin{proof}
\sketch
\leanok
\uses{lem:forward-difference-quotient}
On the compact image of
$\mathcal Z\times[-\eta/2,\eta/2]$, bound $|\chi F|$ by $C$.
Transporting maximizers between nearby fibres shows
$|F_{\max}(t_2)-F_{\max}(t_1)|\le C|t_2-t_1|$. If a forward quotient
at $t_0$ exceeded $c=\psi(t_0,F_{\max}(t_0))$, flow a maximizer at a
nearby later time backward to the $t_0$-fibre. The mean value theorem
and compactness would yield a maximizer $x_0$ at $t_0$ with
$\chi F(x_0)>c$, contradicting the hypothesis. Apply the comparison
lemma to conclude.
\end{proof}

\begin{proposition}[Forward difference maximum property]
\label{prop:forward-difference-maximum}
\lean{MorganTianLib.levelMax_le_of_dir_le}
\source{morgan-tian}
\uses{def:forward-difference-quotient}
\leanok
\label{fordiffmax}




Let $M$ be smooth, let $\chi$ be a smooth vector field, and let
${\bf t}\colon M\to[a,b]$ be smooth with $\chi({\bf t})=1$. Suppose
$F\colon M\to\Ar$ is smooth, attains its maximum on every level of
${\bf t}$, and the set $\mathcal Z$ of all levelwise maximizers is
compact. If
$$\chi(F)(x)\le\psi({\bf t}(x),F(x))\qquad(x\in\mathcal Z),$$
where $\psi\in C^1([a,b]\times\Ar)$, then
$F_{\max}(t)=\max_{{\bf t}^{-1}(t)}F$ is continuous and satisfies
$$\frac{dF_{\max}}{dt}(t)\le\psi(t,F_{\max}(t))$$
in the forward-difference sense. If $G'=\psi(t,G)$ and
$F_{\max}(a)\le G(a)$, then $F_{\max}\le G$ on $[a,b]$.
\end{proposition}
\begin{proof}
\sketch
\leanok
\uses{lem:vector-field-flow-near-compact,lem:forward-difference-maximum-flow}
The local flow of $\chi$ near $\mathcal Z$ moves ${\bf t}$ at unit
speed and differentiates $F$ by $\chi(F)$. Therefore all hypotheses of
Lemma~\ref{lem:forward-difference-maximum-flow} hold, and its conclusions
are exactly the assertions.
\end{proof}
